\section{Disjoint interiors of the initial regions}
\label{sec:area_law_initial_region_interiors}

The initial regions in \cite[Section~11]{OpenAI2026AreaLaw} may share
boundary segments or isolated points. Their open interiors are disjoint.

% Source: 10-geometry.tex, prop:two-families, lines 154--177,
% 212--218 and 299--323, revision adc7f1241b42e322a6451854ab7e4b4c146bf78a.
% These assertions concern the initial family before repairs.

\begin{theorem}[Distinct runs have disjoint interiors]
    \label{thm:area_law_fan_run_disjoint_interiors}
    \lean{TNLean.PEPS.AreaLaw.Geometry.cellFanRunRegions_disjoint_interiors}
    \leanok
    \uses{def:area_law_fan_runs}
    Fix any translated dyadic square, optional midpoint mask and assignment
    of two colors to its fan triangles. If $R\ne T$ are two runs, then
    $\operatorname{int}(T_R)\cap\operatorname{int}(T_T)=\varnothing$.
\end{theorem}
\begin{proof}\leanok
    \uses{thm:area_law_fan_triangle_contact,thm:area_law_cell_fan_intersections,
        def:area_law_fan_runs}
    A real segment in the plane is nowhere dense: it is closed, and the
    affine span of its two endpoints is proper, so its interior is empty.
    Distinct runs have disjoint sets of constituent triangles. For every
    triangle $P_e$ in $R$ and $P_f$ in $T$, their intersection is either
    an actual radial segment or the common center. In either case it is
    nowhere dense. Since both runs have finitely many triangles,
    \begin{align}
        T_R\cap T_T
         & \subseteq\bigcup_{e\in R}\bigcup_{f\in T}(P_e\cap P_f)
    \end{align}
    is nowhere dense. Hence its interior, which equals
    $\operatorname{int}(T_R)\cap\operatorname{int}(T_T)$, is empty.
\end{proof}

\begin{theorem}[Pairwise disjoint initial open regions]
    \label{thm:area_law_initial_open_disjoint}
    \lean{TNLean.PEPS.AreaLaw.Geometry.initialOpenRegion_pairwise_disjoint}
    \leanok
    \uses{def:area_law_initial_open_regions}
    Fix an arbitrary origin, finite endpoint set, integer width $C\ge 2$
    and initial layer $k_0\ge 50{,}000{,}000$, with arbitrary valid
    residue choices at every layer. For any distinct actual initial
    identifiers $i,j$,
    $\operatorname{int}(X_i^0)\cap\operatorname{int}(X_j^0)=\varnothing$.
    No nonemptiness condition on the endpoint set is required.
\end{theorem}
\begin{proof}\leanok
    \uses{thm:area_law_fan_run_disjoint_interiors,
        thm:area_law_primary_fine_cell_cover,thm:area_law_fine_cells_disjoint,
        thm:area_law_dummy_layer_disjoint,thm:area_law_fan_run_cover,
        def:area_law_initial_birth_regions}
    A half-open dyadic square $Q$ is convex with nonempty interior, so
    $\overline{\operatorname{int}(Q)}=\overline Q$. Suppose that $Q$
    is disjoint from each half-open cell $T_i$ in an opposing family.
    Since $\operatorname{int}(Q)$ is open, it is also disjoint from each
    closure $\overline{T_i}$, and thus from their union. The open set
    $U=\operatorname{int}\bigl(\bigcup_i\overline{T_i}\bigr)$ is
    therefore disjoint from $\operatorname{int}(Q)$, hence from its
    closure:
    \begin{align}
        \overline Q\cap\operatorname{int}\Bigl(\bigcup_i\overline{T_i}\Bigr)
         & =\varnothing.
    \end{align}
    Union over the first family's cells to obtain disjoint interiors
    of the two unions of closures.

    Each primary is exactly a finite union of closed nonbelt fine cells,
    while every run is contained in its own closed cell. Different
    primary identifiers cannot share an indexed nonbelt cell: a cell
    has one layer and one primary pitch index. A belt cell cannot be
    such a nonbelt cell. Distinct indexed actual fine cells are disjoint
    before closure, so the preceding closure argument treats all
    different-cell cases. Later cells are likewise disjoint from the
    half-open dummy neighborhood; the same two open-set closure steps
    treat dummy contacts. The only remaining case is two distinct runs
    in one cell, for which the first theorem applies.
\end{proof}
